% Part: propositional-logic
% Chapter: propositional-logic
% Section: preliminaries

\documentclass[../../../include/open-logic-section]{subfiles}

\begin{document}

\olfileid[fr]{pl}{syn}{pre}

\olsection{Préliminaires}

\begin{thm}[\emph{Principe d'induction sur les !!{formula}s}]
  \ollabel{thm:induction}
  Soit $P$ une propriété vérifiée par toutes les !!{formula}s atomiques
  et telle que :
  \begin{enumerate}
    \tagitem{prvNot}{si elle est vérifiée par~$!A$, elle l'est
      par $\lnot !A$ ;}{}
    \tagitem{prvAnd}{si elle est vérifiée par $!A$ et~$!B$,
      elle l'est par $(!A \land !B)$ ;}{}
    \tagitem{prvOr}{si elle est vérifiée par $!A$ et~$!B$,
      elle l'est par $(!A \lor !B)$ ;}{}
    \tagitem{prvIf}{si elle est vérifiée par $!A$ et~$!B$,
      elle l'est par $(!A \lif !B)$ ;}{}
    \tagitem{prvIff}{si elle est vérifiée par $!A$ et~$!B$,
      elle l'est par $(!A \liff !B)$ ;}{}
  \end{enumerate}
  Alors $P$ est vérifiée par toutes les !!{formula}s.
\end{thm}

\begin{proof}
  Soit $S$ la collection de toutes les !!{formula}s qui vérifient~$P$.
  On a évidemment $S \subseteq \Frm[L_0]$. De plus, $S$ remplit toutes
  les conditions de la \olref[fml]{defn:formulas} : il contient toutes
  les !!{formula}s atomiques et il est stable par les !!{operator}s.
  Comme $\Frm[L_0]$ est la plus petite classe ayant ces propriétés,
  $\Frm[L_0] \subseteq S$. Ainsi $\Frm[L_0] = S$, et toute formule
  vérifie~$P$.
\end{proof}

\begin{prop}\ollabel{prop:balanced}
   Toute !!{formula} de~$\Frm[L_0]$ est \emph{équilibrée} :
   elle comporte autant de parenthèses ouvrantes que de parenthèses
   fermantes.
\end{prop}

\begin{prob}
Démontrer la \olref[pl][syn][pre]{prop:balanced}.
\end{prob}

\begin{prop}\ollabel{prop:noinit}
Aucun préfixe propre d'une !!{formula} n'est une !!{formula}.
\end{prop}

\begin{prob}
  Démontrer la \olref[pl][syn][pre]{prop:noinit}.
\end{prob}

\begin{prop}[Lecture unique]
Toute !!{formula}~$!A$ de $\Frm[L_0]$ se décompose selon exactement
l'un des cas suivants :
\begin{enumerate}
\tagitem{prvFalse}{$\lfalse$.}{}

\tagitem{prvTrue}{$\ltrue$.}{}

\item $\Obj p_n$, pour un certain $\Obj p_n \in \PVar$.

\tagitem{prvNot}{$\lnot !B$, pour une !!{formula}~$!B$.}{}

\tagitem{prvAnd}{$(!B \land !C)$, pour des !!{formula}s $!B$ et~$!C$.}{}

\tagitem{prvOr}{$(!B \lor !C)$, pour des !!{formula}s $!B$ et~$!C$.}{}

\tagitem{prvIf}{$(!B \lif !C)$, pour des !!{formula}s $!B$ et~$!C$.}{}

\tagitem{prvIff}{$(!B \liff !C)$, pour des !!{formula}s $!B$ et~$!C$.}{}
\end{enumerate}
De plus, cette décomposition est \emph{unique}.
\end{prop}

\begin{proof}
Par induction sur $!A$. Supposons par exemple que $!A$ admette deux
lectures distinctes, $(!B \lif !C)$ et $(!B' \lif !C')$. Alors $!B$
et $!B'$ doivent être identiques, sinon l'une serait un préfixe propre
de l'autre. Si les deux lectures de $!A$ sont distinctes, $!C$ et $!C'$
doivent donc être deux lectures distinctes de la même chaîne de
symboles, ce qui contredit l'hypothèse d'induction.
\end{proof}

\begin{defn}[Substitution uniforme]
Si $!A$ et $!B$ sont des !!{formula}s et $\Obj p_i$ une
!!{propositional variable}, on note $\Subst{!A}{!B}{\Obj p_i}$ le
résultat du remplacement de chaque occurrence de $\Obj p_i$ dans $!A$
par une occurrence de $!B$. De même, on note
$\SSubst{!A}{\subst{!B_1}{\Obj p_1},\dots,\subst{!B_n}{\Obj p_n}}$
la substitution simultanée des !!{formula}s $!B_1$, \dots,~$!B_n$
aux variables $\Obj p_1$, \dots,~$\Obj p_n$.
\end{defn}

\begin{prob}
Pour chacune des cinq !!{formula}s ci-dessous, déterminer si elle
peut s'écrire comme une substitution \(\Subst{!A}{!B}{\Obj p_i}\)
où \(!A\) est (i)~\(\Obj p_0\) ; (ii)~\((\lnot \Obj p_0 \land \Obj
p_1)\) ; (iii)~\(((\lnot \Obj p_0 \lif \Obj p_1) \land \Obj
p_2)\). Préciser dans chaque cas la substitution correspondante.
  \begin{enumerate}
    \item \(\Obj p_1\)
    \item \((\lnot \Obj p_0 \land \Obj p_0)\)
    \item \(((\Obj p_0 \lor \Obj p_1) \land \Obj p_2)\)
    \item \(\lnot ((\Obj p_0 \lif \Obj p_1) \land \Obj p_2)\)
    \item \(((\lnot (\Obj p_0 \lif \Obj p_1) \lif (\Obj p_0 \lor \Obj p_1)) \land \lnot (\Obj p_0 \land \Obj p_1))\)
  \end{enumerate}
\end{prob}

\begin{prob}
Donner une définition mathématiquement rigoureuse de $\Subst{!A}{!B}{p}$
par induction.
\end{prob}

\end{document}
